\section{Parameter provenance and independent estimation}
\label{sec:estimation}

\subsection{Provenance}

Each parameter is classified as \emph{calibrated} (fixed by market prices),
\emph{inherited} (taken from a historical estimation under $\PP$ without
adjustment), \emph{derived} (reconstructed from summary statistics of that
estimation, without a standard error) or \emph{assumed}. Only the forward curve
is calibrated (Table~\ref{tab:provenance}). The transition intercepts were not
reported by the original estimation and were reconstructed by matching the
reported mean regime durations of $3.76$ and $7.56$ hours
(\ref{app:transitions}). The scale $s_P$ is documented as a training
median of absolute prices, but the window that produced $282.48$\,\TRY{} is not
available: the same definition gives $1{,}399.99$\,\TRY{} on 2019--2025 and
$302.02$\,\TRY{} on 2019--2020, so the value is consistent with an early,
pre-depreciation window.

\begin{table}[htbp]
\centering
\caption{Production parameters, their provenance, and maximum-likelihood
estimates on the model-consistent residual over 2019--2025 (standard errors in
parentheses). Only $F(\cdot)$ is identified by market prices.}
\label{tab:provenance}
\small
\setlength{\tabcolsep}{4pt}
\begin{tabular}{@{}lllr@{}}
\toprule
Parameter & Production & Status & Re-estimated \\
\midrule
$F(\cdot)$ & exact on six quotes & Calibrated & -- \\
January level & linear anchor & Assumed & -- \\
$\kappa$ (h$^{-1}$) & $0.0784$ & Inherited & $0.152$ \\
$\sigma^y_0$ & $0.00354$ & Inherited & $0.078\ (0.006)$ \\
$\sigma^y_1$ & $0.0924$ & Inherited & $0.320\ (0.011)$ \\
$\alpha_{01}$ & $-1.016$ & Derived & $-2.615\ (0.038)$ \\
$\alpha_{10}$ & $-1.895$ & Derived & $-1.071\ (0.038)$ \\
$\gamma_{01}$ & $-0.584$ & Inherited & $-0.759\ (0.034)$ \\
$\gamma_{10}$ & $+0.078$ & Inherited & $+0.692\ (0.041)$ \\
$s_P$ (\TRY) & $282.48$ & Assumed & -- \\
$\pi_1$ at last known hour & $0.068$ & Inherited & -- \\
$a_i$, $\eta_{ij}$ & $0$ & Assumed & -- \\
$r$ (annual) & $0.40$ & Assumed & -- \\
\bottomrule
\end{tabular}
\end{table}

\subsection{Re-estimation on a model-consistent residual}

The inherited parameters were tested by estimating the residual process
independently. The object that \eqref{eq:ou} describes is a deviation from the
forward curve, so the estimation series is built to match it: the transformed
price minus its calendar-month mean minus the within-month hour-of-week shape.
The two-regime model with time-varying transitions is estimated by maximum
likelihood with the filter of \citet{Hamilton1989}, using multiple starting
points and Hessian-based standard errors (\ref{app:estimation}).

Three results stand out. First, mean reversion on this residual is fast:
single-regime fits give $\kappa$ between $0.19$ and $0.26$\,h$^{-1}$ in every
year from 2019 to 2025, and the two-regime fit gives $\kappa=0.152$\,h$^{-1}$
with an interior autoregressive coefficient of $0.859$. The production value
of $0.078$\,h$^{-1}$ is reached in no year. Second, the production innovation
scale is too narrow, $222$ against $359$\,\TRY{} per hour in 2025. Third,
time-varying transitions are strongly supported: the likelihood-ratio
statistic against constant transitions is $1{,}851$ on two degrees of freedom.

The origin of the inherited persistence is visible from a fit to the
transformed price level itself, without the monthly level removed. That fit
places the autoregressive coefficient on the unit-root boundary, as the
inherited estimate does, because over 2019--2025 the level is dominated by
lira depreciation. A two-by-two comparison of scale and level removal
confirms the attribution: removing the calendar-month level changes the
estimated $\kappa$ by a factor of six to ten, whereas changing the scale from
the transform to lira changes it by a factor of 1.3 to 1.7. An earlier version
of the present model inherited $\kappa=4.11\times10^{-6}$\,h$^{-1}$, a
half-life of nineteen years, from such a fit and priced the reference call at
$687.04$\,\TRY, four times its present value. Persistence of the transformed
price level and persistence of deviations around a forward curve are different
objects, and only the second belongs in \eqref{eq:ou}.

\subsection{What the inherited parameters get right}
\label{sec:statvar}

Mean reversion and innovation scale are individually wrong in opposite
directions, and the errors offset in the quantity that governs option values
at the reported maturities, the stationary variance. A regime-mixture
approximation gives a stationary residual standard deviation of
$582.5$\,\TRY{} at the valuation forward, and along the covariate path used in
pricing the terminal standard deviation at 72 hours is $532.3$\,\TRY. The
realised value for 2025 is $531$\,\TRY{} when the hour-of-week shape is
estimated on 2025 alone and $624$\,\TRY{} when it is pooled over 2019--2025.
In a long-path simulation the Kolmogorov--Smirnov distance to the realised 2025
residual is $0.087$ for the production parameters, against $0.125$ for the
full-window re-estimate and $0.159$ for a fit restricted to 2022--2025; no
parameter set is accepted by the test, partly because the realised series is
truncated by the price limits. The re-estimated sets imply stationary
deviations of $929$ and $1{,}088$\,\TRY{} and overstate 2025 dispersion.
The production parameters are therefore retained for pricing, and the
re-estimated full-window set is carried as an alternative in the out-of-sample
evaluation.

Holding the stationary variance at its production value isolates the role of
$\kappa$ (Figure~\ref{fig:sensitivity}a). In a single-regime reduction, the
72-hour call at $K=3{,}000$ stays between $169$ and $176$\,\TRY{} across the
estimated range $\kappa\in[0.15,0.22]$\,h$^{-1}$ and between $162$ and
$187$\,\TRY{} across $[0.01,0.30]$, while the 1-hour call moves from $7$ to
$98$\,\TRY. Mean reversion therefore matters at short maturities and not at the
24 to 72 hours reported here, and values below about 12 hours are not reliable
under the production parameters. The same reduction also measures what the
regime mixture contributes at equal variance: at the production $\kappa$ the
single-regime call is $183.4$\,\TRY{} against $166.75$\,\TRY{} for the
two-regime process, a difference of $-9\%$ that reflects the heavier tails of
the mixture.
